% SPDX-License-Identifier: Apache-2.0
\section{An Asynchronous FIFO}
\label{sec:x-fifo}

The producer port is in one clock and the consumer port in another,
and the two stand in a stated relation: the write clock has period 2,
the read clock period 3 and phase 1, in the unit the design shares.
Each pointer lives in its own domain and reaches the other side through
a \code{Crossing} that the reading process steps at its own edge, which
is what a synchroniser does. The memory is written by the producer
clock and read by the consumer clock without a wait, as a dual-port
RAM's read port is; the pointer discipline is what makes that safe.
The channel's own handshake is the backpressure: the producer offers
only when the channel is ready, the FIFO takes only when not full, and
the trace shows the producer stalled on \emph{full} twice in twenty
steps.

\lstinputlisting[caption={\code{ex\_fifo.rs}. Run with \code{bazel run //lib/examples:ex\_fifo}.},label={lst:x-fifo}]{lib/examples/ex_fifo.rs}

\lstinputlisting[style=ascii,caption={What \code{ex\_fifo} printed when the build ran it.},label={lst:x-fifo-out}]{out_fifo.txt}
